\chapter{Optimization with Gradients}

Where is the highest point on a mountain? Where does a company make the most profit, or a machine-learning model
make the fewest mistakes? Each question asks for the maximum or minimum of a function of several variables. In
the Gradients chapter, you learned that the gradient points uphill. This chapter uses that fact to find peaks and
valleys, both exactly with algebra and step by step with gradient descent, and then meets some of the famous
equations of physics and economics that are written with partial derivatives.

\section{Optimization with Gradients}\index{optimization!multivariable}

At the top of a hill, there is no uphill direction. The gradient is the
zero vector. The same is true at the bottom of a valley. So to find the
highest and lowest points of $f(x, y)$, we look for \newterm{critical
points}\index{critical point!multivariable}, where
$$\nabla f = [0, 0], \quad\text{that is,}\quad f_x = 0 \text{ and } f_y = 0$$

In one variable, a critical point could be a max, a min, or neither. In two
variables there is a new possibility: a \newterm{saddle
point}\index{saddle point}. The surface $f(x, y) = x^2 - y^2$ at the origin
is shaped like a horse's saddle (or a Pringles chip): walking east-west, the
origin is a low point; walking north-south, it is a high point. Its gradient
is zero, but it is neither a max nor a min.

\begin{mdframed}[style = important, frametitle = {The Second Derivative Test for Two Variables}]
Suppose $\nabla f(a, b) = [0, 0]$. Let
$$D = f_{xx}(a, b)\, f_{yy}(a, b) - \left[f_{xy}(a, b)\right]^2$$
\begin{itemize}
\item If $D > 0$ and $f_{xx}(a, b) > 0$, then $f(a, b)$ is a local minimum.
\item If $D > 0$ and $f_{xx}(a, b) < 0$, then $f(a, b)$ is a local maximum.
\item If $D < 0$, then $(a, b)$ is a saddle point.
\item If $D = 0$, the test tells you nothing.
\end{itemize}
\end{mdframed}

For $x^2 - y^2$: $f_{xx} = 2$, $f_{yy} = -2$, $f_{xy} = 0$, so $D = -4 < 0$:
a saddle, as promised.

\textbf{Example}: Find and classify the critical points of $f(x, y) = x^2 + xy +
y^2 - 3x$.

\textbf{Solution}: Set both partial derivatives to zero:
$$f_x = 2x + y - 3 = 0 \qquad f_y = x + 2y = 0$$
From the second equation, $x = -2y$. Substituting into the first, $-4y + y
- 3 = 0$, so $y = -1$ and $x = 2$. The second partials are $f_{xx} = 2$,
$f_{yy} = 2$, and $f_{xy} = 1$, so $D = (2)(2) - 1^2 = 3 > 0$. Since
$f_{xx} > 0$, the point $(2, -1)$ is a local minimum, with $f(2, -1) =
-3$.

\subsection{Gradient Descent}\index{gradient descent}

Solving $\nabla f = [0, 0]$ by algebra only works for simple functions. A
computer can instead do what a lost hiker in the fog would do: feel which
way is downhill and take a small step that way. Over and over. Because
$\nabla f$ points uphill, each step goes in the direction of $-\nabla f$:
$$[x, y] \leftarrow [x, y] - r \nabla f(x, y)$$
where $r$ is a small step size. This algorithm, \newterm{gradient
descent}, is how most machine learning models are trained: the
``function'' is the model's error, with millions of variables.

\begin{minted}{python}
x, y = 0.0, 0.0     # starting guess
r = 0.1             # step size

for step in range(100):
    fx = 2*x + y - 3    # partial derivatives of x^2 + xy + y^2 - 3x
    fy = x + 2*y
    x = x - r * fx
    y = y - r * fy

print(x, y)
\end{minted}

After 100 steps, this prints values within 0.0001 of $(2, -1)$, the
minimum we found by hand.

\begin{Exercise}[title = {Finding a Maximum}, label = grad_max]
A company sells two products. If it charges $x$ dollars more than a base
price for the first and $y$ dollars more for the second, its profit (in
thousands of dollars) is
$$P(x, y) = 6x + 4y - x^2 - y^2 - xy$$
Find the critical point, show that it is a maximum, and find the maximum
profit.
\end{Exercise}
\begin{Answer}[ref = grad_max]
$P_x = 6 - 2x - y = 0$ and $P_y = 4 - 2y - x = 0$. From the second, $x = 4
- 2y$; substituting, $6 - 8 + 4y - y = 0$, so $y = \frac{2}{3}$ and $x =
\frac{8}{3}$. $P_{xx} = -2$, $P_{yy} = -2$, $P_{xy} = -1$, so $D = 4 - 1 =
3 > 0$ with $P_{xx} < 0$: a maximum. $P\left(\frac{8}{3},
\frac{2}{3}\right) = \frac{28}{3} \approx 9.33$, or about \$9,330.
\end{Answer}

\begin{Exercise}[title = {Saddle or Not?}, label = grad_saddle]
Find and classify the critical point of $f(x, y) = xy$.
\end{Exercise}
\begin{Answer}[ref = grad_saddle]
$f_x = y$ and $f_y = x$, so the only critical point is $(0, 0)$. $f_{xx} =
0$, $f_{yy} = 0$, $f_{xy} = 1$, so $D = 0 - 1 = -1 < 0$: a saddle point.
(Along $y = x$, $f = x^2$ rises; along $y = -x$, $f = -x^2$ falls.)
\end{Answer}

\section{Famous Equations with Partial Derivatives}

\subsection{Laplace's Equation}
A partial differential equation that has applications in fluid dynamics and 
electronics is Laplace's Equation\index{Laplace's equation}. Solutions to 
Laplace's Equation are called \index{harmonic function}\textit{harmonic functions}. 

\begin{mdframed}[style = important, frametitle = {Laplace's Equation}]
Consider a twice-differentiable function, $f$. In two dimensions, Laplace's 
Equation is given by:
$$ \frac{\partial ^2 f}{\partial x^2} + \frac{\partial^2 f}{\partial y^2} = 0$$
And in three dimensions, 
$$\frac{\partial ^2 f}{\partial x^2} + \frac{\partial ^2 f}{\partial y^2} + 
\frac{\partial ^2 f}{\partial z^2} = 0$$
\end{mdframed}

Another way to represent Laplace's Equation is:
$$\Delta f = \nabla^2 f = \nabla \cdot \nabla f = 0$$

Where $\nabla ^2 = \Delta$ is called the \textit{Laplace operator}. 

\textbf{Example}: Determine whether or not $f = x^2 + y^2$ is a solution to 
Laplace's Equation. 

\textbf{Solution}: We are checking to see if $\frac{\partial^2 f}{\partial 
x^2} + \frac{\partial^2 f}{\partial y^2} = 0$ for $f(x, y) = x^2 + y^2$. 
Finding $\partial^2 f / \partial x^2$:
$$\frac{\partial^2 f}{\partial x^2} = \frac{\partial}{\partial x} \left[ \frac{
\partial}{\partial x} \left( x^2 + y^2 \right) \right]$$
$$= \frac{\partial}{\partial x} \left( 2x \right) = 2$$

And finding $\partial^2 f / \partial y^2$:
$$\frac{\partial^2 f}{\partial y^2} = \frac{\partial}{\partial y} \left[ \frac{
\partial}{\partial y} \left( x^2 + y^2 \right) \right]$$
$$= \frac{\partial}{\partial y} \left( 2y \right) = 2$$

Then $\frac{\partial^2 f}{\partial x^2} + \frac{\partial^2 f}{\partial y^2} = 
2 + 2 = 4 \neq 0$. Therefore, $f(x, y) = x^2 + y^2$ is not a solution to 
Laplace's Equation. 

\begin{Exercise}[title = {Solutions to Laplace's Equation}, label = laplace]
Determine whether the function is a solution to Laplace's Equation. 
\begin{enumerate}
    \item $f(x, y) = x^2 - y^2$
    \item $f(x, y) = \sin{x} \cosh{y} + \cos{x} \sinh{y}$
    \item $f(x, y) = e^{-x} \cos{y} - e^{-y} \cos{x}$
\end{enumerate}
\end{Exercise}

\begin{Answer}[ref = laplace]
\begin{enumerate}
    \item $$\frac{\partial^2 f}{\partial x^2} + \frac{\partial^2 f}{\partial 
    y^2} = \frac{\partial}{\partial x} \left[ \frac{\partial}{\partial x} 
    \left( x^2 - y^2 \right) \right] + \frac{\partial}{\partial y} \left[ 
    \frac{\partial}{\partial y} \left( x^2 - y^2 \right) \right]$$
    $$= \frac{\partial}{\partial x} \left[ 2x \right] + \frac{\partial}{
    \partial y} \left[ -2y \right] = 2 - 2 = 0$$

    Therefore, $f(x, y) = x^2 - y^2$ is a solution to Laplace's Equation. 

    \item $$\frac{\partial^2}{\partial x^2} + \frac{\partial^2}{\partial y^2} 
    = \frac{\partial}{\partial x} \left[ \frac{\partial}{\partial x} \left( 
    \sin{x} \cosh{y} + \cos{x} \sinh{y} \right) \right] + \frac{\partial}{
    \partial y} \left[ \frac{\partial}{\partial y} \left( \sin{x} \cosh{y} + 
    \cos{x} \sinh{y} \right) \right]$$
    $$= \frac{\partial}{\partial x} \left[ \cos{x} \cosh{y} - \sin{x} \sinh{y} 
    \right] + \frac{\partial}{\partial y} \left[ \sin{x} \sinh{y} + \cos{x} 
    \cosh{y} \right]$$
    $$= -\sin{x}\cosh{y} - \cos{x}\sinh{y} + \sin{x}\cosh{y} + \cos{x} \sinh{
    y} = 0$$

    Therefore, $f(x, y) = \sin{x}\cosh{y} + \cos{x}\sinh{y}$ is a solution to 
    Laplace's Equation.

    \item $$\frac{\partial^2 f}{\partial x^2} + \frac{\partial^2 f}{\partial 
    y^2} = \frac{\partial}{\partial x} \left[ \frac{\partial}{\partial x} 
    \left( e^{-x}\cos{y} - e^{-y}\cos{x} \right) \right] + \frac{\partial}{
    \partial y} \left[ \frac{\partial}{\partial y} \left( e^{-x} \cos{y} - 
    e^{-y} \cos{x} \right) \right]$$
    $$= \frac{\partial}{\partial x} \left[ -e^{-x}\cos{y} + e^{-y}\sin{x} 
    \right] + \frac{\partial}{\partial y} \left[ -e^{-x}\sin{y} + e^{-y} 
    \cos{x} \right]$$
    $$= e^{-x} \cos{y} + e^{-y} \cos{x} - e^{-x} \cos{y} - e^{-y} \cos{x} = 0$$

    Therefore, $f(x, y) = e^{-x}\cos{y} - e^{-y}\cos{x}$ is a solution to Laplace's Equation.
\end{enumerate}
\end{Answer}

\subsection{The Wave Equation}\index{wave equation}
Another useful equation with partial derivatives is the Wave Equation:
$$\frac{\partial^2 f}{\partial t^2} = a^2 \frac{\partial^2 f}{\partial x^2}$$

Where $f$ is a function of $x$ and $t$ and $a$ is a constant. This equation 
describes waves, such as a vibrating string, light waves, or sound waves. 

\textbf{Example}: Show that $f(x, t) = \sin{ \left( x - at \right)}$ satisfies 
the Wave Equation. 

\textbf{Solution}: First, we find the second partial derivatives:
$$\frac{\partial^2 f}{\partial t^2} = \frac{\partial}{\partial t} \left[ \frac{
\partial}{\partial t} \left( \sin{ \left( x - at \right)} \right) \right] = 
\frac{\partial}{\partial t} \left[ -a\cos{\left(x - at \right)} \right] = -a^2
\sin{\left(x - at \right)}$$

$$\frac{\partial^2 f}{\partial x^2} = \frac{\partial}{\partial x} \left[ \frac{
\partial}{\partial x} \left( \sin{ \left( x - at \right)} \right) \right] = 
\frac{\partial}{\partial x} \left[ \cos{\left( x - at \right)}\right] = -\sin{
\left( x - at \right)}$$

And we see that:
$$a^2 \frac{\partial^2 f}{\partial x^2} = -a^2 \sin{\left(x - at \right)} = 
\frac{\partial^2 f}{\partial t^2}$$

Therefore, this function satisfies the Wave Equation. 

\begin{Exercise}[title = {The Wave Equation}, label = wave]
Show that the following functions satisfy the Wave Equation:
\begin{enumerate}
    \item $f(x, t) = \cos{\left(kx \right)}\cos{\left(akt \right)}$
    \item $f(x, t) = \sin{\left( x - at \right)} + \ln{ \left( x + at 
    \right)}$
    \item $f(x, t) = \frac{t}{a^2t^2 - x^2}$
\end{enumerate}
\end{Exercise}

\begin{Answer}[ref = wave]
\begin{enumerate}
    \item Finding the partial derivatives:
    $$\frac{\partial^2 f}{\partial t^2} = \frac{\partial}{\partial t} \left[ 
    \frac{\partial}{\partial t} \left( \cos{ \left( k x \right)} \cos{ \left( 
    a k t \right)} \right) \right] = \frac{\partial}{\partial t} \left[ -a k 
    \cos{ \left( k x \right)} \sin{ \left( a k t \right)} \right]$$
    $$\frac{\partial^2 f}{\partial t^2} = -a^2 k^2 \cos{ \left( k x \right)} 
    \cos{ \left( a k t \right)}$$

    $$\frac{\partial^2 f}{\partial x^2} = \frac{\partial}{\partial x} \left[ 
    \frac{\partial}{\partial x} \left( \cos{\left( k x \right)}\cos{\left( a k 
    t \right)} \right) \right] = \frac{\partial}{\partial x} \left[ -k \sin{ 
    \left( k x \right)} \cos{ \left( a k t \right)} \right]$$
    $$\frac{\partial^2 f}{\partial x^2} = -k^2 \cos{ \left( k x \right)} \cos{ 
    \left( a k t \right)}$$

    And we see that:
    $$a^2 \frac{\partial^2 f}{\partial x^2} = -a^2 k^2 \cos{ \left( k x 
    \right)} \cos{ \left( a k t \right)} = \frac{\partial^2 f}{\partial t^2}$$

    Therefore, $f(x, t) = \cos{\left( k x \right)}\cos{\left( a k t \right)}$ 
    satisfies the Wave Equation. 

    \item Finding the partial derivatives:
    $$\frac{\partial^2 f}{\partial t^2} = \frac{\partial}{\partial t} \left[ 
    \frac{\partial}{\partial t} \left( \sin{\left( x - a t \right)} + \ln{ 
    \left( x + a t \right)} \right) \right]$$
    $$\frac{\partial^2 f}{\partial t^2} = \frac{\partial}{\partial t} \left[ 
    -a \cos{ \left( x - a t \right)} + \frac{a}{x + a t} \right] = -a^2 \sin{ 
    \left( x - a t \right)} + \frac{-a^2}{\left( x + a t \right)^2}$$

    $$\frac{\partial^2 f}{\partial x^2} = \frac{\partial}{\partial x} \left[ 
    \frac{\partial}{\partial x} \left( \sin{ \left( x - a t \right)} + \ln{ 
    \left( x + a t \right)} \right) \right]$$

    $$\frac{\partial^2 f}{\partial x^2} = \frac{\partial}{\partial x} \left[ 
    \cos{ \left( x - a t \right)} + \frac{1}{x + a t} \right] = -\sin{ \left( 
    x - a t \right)} + \frac{-1}{\left( x + a t \right)^2}$$

    And we see that:
    $$a^2 \frac{\partial^2 f}{\partial x^2} = -a^2 \sin{ \left( x - a t 
    \right)} + \frac{-a^2}{\left( x + a t \right)^2} = \frac{\partial^2 f}{
    \partial t^2}$$

    Therefore, $f(x, t) = \sin{ \left( x - a t \right)} + \ln{ \left( x + a t 
    \right)}$ satisfies the Wave Equation. 

    \item Finding the partial derivatives:
    $$\frac{\partial^2 f}{\partial t^2} = \frac{\partial}{\partial t} \left[ 
    \frac{\partial}{\partial t} \left( \frac{t}{a^2 t^2 - x^2} \right) \right] 
    = \frac{\partial}{\partial t} \left[ \frac{\left(a^2 t^2 - x^2 \right) - t 
    \left(2 a^2 t \right)}{\left( a^2 t^2 - x^2 \right)^2} \right]$$
    $$\frac{\partial^2 f}{\partial t^2} = \frac{\partial}{\partial t} \left[ 
    \frac{a^2 t^2 - x^2 - 2 a^2 t^2}{\left( a^2 t^2 - x^2 \right)^2} \right] = 
    \frac{\partial}{\partial t} \left[ \frac{-a^2 t^2 - x^2}{\left( a^2 t^2 - 
    x^2 \right)^2} \right]$$
    $$\frac{\partial^2 f}{\partial t^2} = \frac{\left( a^2 t^2 - x^2 \right)^2 
    \left(-2 a^2t \right) - \left( -a^2 t^2 - x^2 \right) \left( 2 \left( a^2 
    t^2 - x^2 \right) \left( 2 a^2 t \right) \right)}{\left( a^2 t^2 - x^2 
    \right)^4}$$
    $$\frac{\partial^2 f}{\partial t^2} = \frac{\left( a^2 t^2 - x^2 \right) 
    \left( -2 a^2 t \right) - \left( -a^2 t^2 - x^2 \right) \left( 4 a^2 t 
    \right)}{\left( a^2 t^2 - x^2 \right)^3}$$
    $$\frac{\partial^2 f}{\partial t^2} = \frac{ -2 a^4 t^3 + 2 a^2 t x^2 + 4 
    a^4 t^3 + 4 a^2  t x^2}{\left( a^2 t^2 - x^2 \right)^3}$$
    $$\frac{\partial^2 f}{\partial t^2} = \frac{ 2 a^4 t^3 + 6 a^2 t x^2}{
    \left( a^2 t^2 - x^2 \right)^3} = 2 a^2 t \left( \frac{a^2 t^2 + 3 x^2}{
    \left( a^2 t^2 - x^2 \right)^3} \right)$$

    $$\frac{\partial^2 f}{\partial x^2} = \frac{\partial}{\partial x} \left[ 
    \frac{\partial}{\partial x} \left( \frac{t}{a^2 t^2 - x^2} \right) \right] 
    = \frac{\partial}{\partial x} \left[ \frac{\left( a^2 t^2 - x^2 \right) 
    \left( 0 \right) - t \left( -2 x \right)}{\left( a^2  t^2 - x^2 \right)^2} 
    \right]$$
    $$\frac{\partial^2 f}{\partial x^2} = \frac{\partial}{\partial x} \left[ 
    \frac{2 t x}{\left( a^2 t^2 - x^2 \right)^2} \right] = \frac{\left( a^2 
    t^2 - x^2 \right)^2 \left( 2 t \right) - \left( 2 t x \right) \left( 2 
    \left( a^2 t^2 - x^2 \right) \left( -2 x \right) \right)}{\left( a^2 t^2 - 
    x^2 \right)^4}$$
    $$\frac{\partial^2 f}{\partial x^2} = \frac{\left( a^2 t^2 - x^2 \right) 
    \left( 2 t \right) - \left( 2 t x \right) \left( 2 \right) \left( -2 x 
    \right)}{\left( a^2 t^2 - x^2 \right)^3}$$
    $$\frac{\partial^2 f}{\partial x^2} = \frac{2 a^2 t^3 - 2 t x^2 + 8 t x^2}{
    \left( a^2 t^2 - x^2 \right)^3} = \frac{2 a^2 t^3 + 6 t x^2}{\left( a^2 
    t^2 - x^2 \right)^3} = 2 t \left( \frac{a^2 t^2 + 3 x^2}{\left( a^2 t^2 - 
    x^2 \right)^3} \right)$$

    And we see that:
    $$a^2 \frac{\partial^2 f}{\partial x^2} = 2 a^2 t \left( \frac{a^2 t^2 + 3 
    x^2}{\left( a^2 t^2 - x^2 \right)^3} \right) = \frac{\partial^2 f}{
    \partial t^2}$$

    Therefore, $f(x, t) = \frac{t}{a^2t^2 - x^2}$ satisfies the Wave Equation. 
\end{enumerate}
\end{Answer}

\subsection{Cobb-Douglas Production Function}\index{Cobb-Douglas production function}
The Cobb-Douglas function describes how production depends on capital and labor, 
as theorized by the economists Charles Cobb and Paul Douglas. Capital 
investments are things like new machinery, expanded factories, or raw 
materials. Labor investments involve hiring more workers or improving 
working conditions to improve work rates. We can describe total production, 
$P$, as a function of labor, $L$, and capital, $K$. Cobb and Douglas posit 
three conditions:
\begin{enumerate}
\item Without either labor or capital, production will cease.
\item The marginal productivity of labor is proportional to the amount of production 
per unit of labor.
\item The marginal productivity of capital is proportional to the amount of 
production per unit of capital.
\end{enumerate}

The marginal productivity of labor is given by the partial derivative, $\partial P 
/ \partial L$ and the production per unit of labor is given by $P / L$. 
Therefore, statement 2 says that:
$$\frac{\partial P}{\partial L} = \alpha \frac{P}{L}$$

where $\alpha$ is some constant. Keeping $K$ constant at $K = K_0$, we have 
the differential equation:
$$ \frac{d P}{d L} = \alpha \frac{P}{L}$$

Solving, we find that:
$$P(L, K_0) = C_1 \left( K_0 \right) L^{\alpha}$$

We make $C_1$ a function of $K_0$ because it could depend on $K_0$. In a 
similar manner to above, we can write statement 3 as a mathematical statement:
$$\frac{\partial P}{\partial K} = \beta \frac{P}{K}$$

where $\beta$ is also a constant. Keeping $L = L_0$ and solving, we see that:
$$P(L_0, K) = C_2(L_0)K^{\beta}$$

again, we assume $C_2$ is a function of the fixed labor, $L_0$. Combining 
these equations, we get:
$$P(L, K) = bL^{\alpha} K^{\beta}$$

where $b$ is a constant independent of capital and labor. Additionally, from 
statement 1, we know that $\alpha > 0$ and $\beta > 0$. What happens if both 
labor and capital are increased by a factor of $n$? Let us examine the effect 
on $P$:
$$P(nL, nK) = b \left( nL \right)^{\alpha} \left( nK \right)^{\beta}$$
$$P(nL, nK) = n^{\alpha + \beta} b L^{\alpha} K^{\beta} = n^{\alpha + \beta} 
P(L, K)$$

Cobb and Douglas noted that if $\alpha + \beta = 1$, then $P(nL, nK) = nP(L, 
K)$, and therefore increasing labor and capital by a factor of $n$ increases 
production by a factor of $n$ as well. Therefore, the Cobb-Douglas equation 
assumes $\alpha + \beta = 1$ and can be written as:
$$P(L, K) = bL^{\alpha}K^{1 - \alpha}$$

\begin{Exercise}[title = {Cobb-Douglas Production Model}, label = cobb]
Cobb and Douglas modeled production in the US from 1900 to 1922 with the 
equation $P(L, K) = 1.01L^{0.75} K^{0.25}$. 
\begin{enumerate}
\item Express the marginal productivity of labor as a function of $L$ and $K$.
\item Express the marginal productivity of capital as a function of $L$ and $K$.
\item In 1916, $L = 382$ and $K = 126$ (compared to initial values of $100$ in 
1900). What is the marginal productivity of labor in 1916? Of capital?
\item Based on your answer to the previous question, would you invest in 
capital or labor if you owned a factory in 1916? Why?
\end{enumerate}
\end{Exercise}

\begin{Answer}[ref = cobb]
\begin{enumerate}
    \item The marginal productivity of labor is given by $\partial P / \partial L$:
    $$\frac{\partial P}{\partial L} = \frac{\partial}{\partial L} \left[ 1.01 
    L^{0.75} K^{0.25} \right] = 0.7575 L^{-0.25} K^{0.25}$$
    \item The marginal productivity of capital is given by $\partial P / \partial 
    K$:
    $$\frac{\partial P}{\partial K} = \frac{\partial}{\partial K} \left[ 1.01 
    L^{0.75} K^{0.25} \right] = 0.2525 L^{0.75} K^{-0.75}$$
    \item Finding the marginal productivity of labor in 1916:
    $$\frac{\partial P}{\partial L} = 0.7575 \left( 382 \right)^{-.25} \left( 
    126 \right)^{0.25} \approx 0.574$$
    And finding the marginal productivity of capital in 1916:
    $$\frac{\partial P}{\partial K} = 0.2525 \left( 382 \right)^{0.75} \left( 
    126 \right)^{-0.75} \approx 0.580$$
    \item Since the marginal productivity of capital is greater, I would invest in 
    capital. This would yield a greater increase in production than the same 
    investment in labor. 
\end{enumerate}
\end{Answer}
